\section{Diagnosis of the mixed cross-gradient ansatz}

\subsection{What does commute}

Let the canonical phase-space coordinates be
\[
(\bm{x},\bm{p})=(x^1,\ldots,x^d,p_1,\ldots,p_d)
\]
on an open domain in \(\R^{2d}\). For a sufficiently smooth function \(f(\bm{x},\bm{p})\), Clairaut's theorem gives
\begin{equation}
\frac{\partial^2 f}{\partial x^i\partial p_j}
=
\frac{\partial^2 f}{\partial p_j\partial x^i}.
\end{equation}
Equivalently,
\begin{equation}
\boxed{
\comm{\partial_{x^i}}{\partial_{p_j}}=0.
}
\label{eq:mixedcommute}
\end{equation}
Equation~\eqref{eq:mixedcommute} is the correct zero identity for independent Cartesian phase-space coordinates.

\subsection{Why the cross product does not vanish}

Only when \(d=3\), after choosing coordinate bases and identifying both gradient triples with ordinary Euclidean vectors, can one formally write
\begin{equation}
\left(\nabla_{\bm{x}}\times\nabla_{\bm{p}}\right)_i
=
\epsilon_{ijk}\,
\partial_{x^j}\partial_{p_k}.
\label{eq:formal-cross}
\end{equation}
Commutation of \(\partial_{x^j}\) with \(\partial_{p_k}\) does not force the antisymmetric contraction over different index pairs to vanish.

\begin{proposition}
The differential operator in Eq.~\eqref{eq:formal-cross} is generally nonzero.
\end{proposition}

\begin{proof}
Take
\[
f(\bm{x},\bm{p})=x^1p_2.
\]
Then
\begin{align}
\left(\nabla_{\bm{x}}\times\nabla_{\bm{p}}\right)_3f
&=
\left(
\partial_{x^1}\partial_{p_2}
-
\partial_{x^2}\partial_{p_1}
\right)x^1p_2
\\
&=1.
\end{align}
Therefore \(\nabla_{\bm{x}}\times\nabla_{\bm{p}}\) is not the zero operator.
\end{proof}

This counterexample also exposes the conceptual error: the equality of two mixed derivatives with the \emph{same} index pair does not imply cancellation between different pairs.

\subsection{Four independent defects in the ansatz}

Equation~\eqref{eq:bad} fails for at least four reasons.

\paragraph{1. It is not an operator identity.}
The left-hand side is generally nonzero, as shown above, and its action depends on the function to which it is applied.

\paragraph{2. It has a tensor-type mismatch.}
The formal cross product is vector-valued, whereas \(\hbar^2/(2\pi)\) is a scalar. A direction or covariant tensorial construction has not been supplied.

\paragraph{3. It is dimensionally inconsistent.}
The mixed differential operator carries dimensions
\begin{equation}
\left[\partial_x\partial_p\right]
=
\frac{1}{LP}
=
\frac{1}{\text{action}},
\end{equation}
while
\begin{equation}
[\hbar^2]=(\text{action})^2.
\end{equation}
Multiplication by a numerical factor such as \(1/(2\pi)\) cannot repair the mismatch.

\paragraph{4. It is not intrinsic to phase space.}
A cross product is a special operation associated with oriented three-dimensional inner-product spaces. A canonical phase space has dimension \(2d\), and the position and momentum derivatives naturally belong to different coordinate directions on a cotangent bundle. No canonical cross product between them exists for general \(d\) or under general canonical coordinate transformations.

The correct conclusion is therefore not that the left-hand side must be replaced by zero, but that no universal scalar equality of the form in Eq.~\eqref{eq:bad} is mathematically available.
